% ===================================================================
% The four resolved overclaims: mathematically true forms
%
% These constructions replace the overclaimed versions with their
% honest, proved forms. Each states precisely what is proved,
% what is conditional, and what is the checkable computation
% that would upgrade a conditional statement to a theorem.
% ===================================================================

% -------------------------------------------------------------------
\subsection{The universal bulk and its specializations}
\label{subsec:universal-bulk-specializations}
% -------------------------------------------------------------------

The Swiss-cheese theorem
(Theorem~\ref{V1-thm:thqg-swiss-cheese})
establishes an abstract universality: for any curved
$A_\infty$-chiral algebra~$\cA$, the chiral derived center
$\cZ^{\mathrm{der}}_{\mathrm{ch}}(\cA)
= H^*(C^\bullet_{\mathrm{ch}}(\cA, \cA), \delta)$
is the universal bulk. The question ``is this the \emph{known}
bulk of a specific physical theory?'' requires additional input.

\begin{proposition}[Specialization to the boundary-linear sector;
 \ClaimStatusProvedHere]
\label{prop:bulk-specialization-boundary-linear}
For a logarithmic $\SCchtop$-algebra with boundary algebra
$A_\partial$ in the boundary-linear exact sector:
\[
 \cO_{\mathrm{bulk}}
 \;\simeq\;
 C^\bullet_{\mathrm{ch}}(A_\partial, A_\partial)
 \;\simeq\;
 \cZ^{\mathrm{der}}_{\mathrm{ch}}(\cC_\partial),
\]
where $\cC_\partial$ is the open-sector dg-category and the
second quasi-isomorphism requires compact generation.
\end{proposition}

\begin{proof}
The first quasi-isomorphism is
Theorem~\ref{thm:bulk_hochschild} (unconditional for any
logarithmic $\SCchtop$-algebra). The second follows from
Keller's identification of the derived center of a compactly
generated dg-category with the Hochschild cochains of any compact
generator's endomorphism algebra~\cite{Kel06}.
\end{proof}

\begin{remark}[The checkable computation]
\label{rem:checkable-bulk-computation}
For Chern--Simons at integrable level~$k$ with gauge group~$G$:
the passage from the chain-level derived center
$\cZ^{\mathrm{der}}_{\mathrm{ch}}(V_k(\fg))$ to the known
Verlinde ring requires:
\begin{enumerate}[label=\textup{(\roman*)}]
\item computing $HH^0(V_k(\fg), V_k(\fg))$, which should equal
 the center $Z(\operatorname{Rep}(V_k(\fg)))$ of the
 module category;
\item at integrable levels, the module category is semisimple
 (Huang--Lepowsky), so its center is the Verlinde algebra
 $\bigoplus_\lambda \C \cdot [\lambda]$.
\end{enumerate}
For $\fg = \mathfrak{sl}_2$ at $k = 1$: the module category has
$2$ simple objects ($L_0, L_1$), so $\dim Z = 2$. The
computation $\dim HH^0(V_1(\mathfrak{sl}_2)) = 2$ would give
the derived-center-as-bulk claim its first checkable confirmation.
This computation is within reach of the existing compute
infrastructure.
\end{remark}

% -------------------------------------------------------------------
\subsection{The line-operator category: MC couplings vs.\ Koszul
 modules}
\label{subsec:line-mc-vs-koszul}
% -------------------------------------------------------------------

\begin{definition}[MC coupling category]
\label{def:mc-coupling-category}
\index{MC coupling category|textbf}
\index{line operators!MC coupling presentation|textbf}
For any $A_\infty$-chiral algebra~$\cA$, the \emph{MC coupling
category} $\operatorname{Line}(\cA)$ has:
\begin{enumerate}[label=\textup{(\alph*)}]
\item \emph{Objects}: pairs $(V, \mu)$ where $V$ is a dg vector
 space and $\mu \in \End(V) \widehat\otimes \cA$ of degree~$1$
 satisfies the $A_\infty$ Maurer--Cartan equation
 $\sum_{k \ge 1} m_k(\mu, \ldots, \mu) = 0$.

\item \emph{Morphisms}:
 $\Hom((V,\mu), (W,\nu))
 = \Hom(V, W) \widehat\otimes \cA$
 with twisted differential
 $d_{\mu,\nu}(f)
 = \sum_{r,s \ge 0} m_{r+s+1}(\nu^{\otimes r}, f,
 \mu^{\otimes s})$.
\end{enumerate}
This category is \emph{always} well-defined: it requires only the
$A_\infty$ structure on~$\cA$ and convergence of the completed
tensor product, not Koszulness or any other structural hypothesis.
It is the intrinsic primitive object.
\end{definition}

\begin{theorem}[Koszul comparison; \ClaimStatusProvedHere]
\label{thm:koszul-comparison-mc}
\index{Koszul comparison!MC coupling|textbf}
If $\cA$ is chirally Koszul
\textup{(}bar cohomology concentrated in bar degree~$1$;
any of the $10$ unconditional equivalent conditions of
Theorem~\textup{\ref{V1-thm:koszul-equivalences-meta}}\textup{)},
then the Koszul dual $\cA^!$ is well-defined and the
comparison functor
\begin{align*}
 \Phi \colon \operatorname{Line}(\cA)
 &\;\longrightarrow\; \Perf(\cA^!) \\
 (V, \mu) &\;\longmapsto\;
 (V,\, \rho = (\mathrm{id}_V \otimes \pi)(\mu))
\end{align*}
is a quasi-equivalence, where
$\pi \colon \cA \twoheadrightarrow \cA^!$ is the bar-cohomology
projection.
\end{theorem}

\begin{proof}
On the Koszul locus, bar-cobar inversion
(Theorem~B) gives a quasi-isomorphism
$\Omega(\barB(\cA)) \xrightarrow{\sim} \cA$. The universal
twisting morphism $\tau \in \cA^! \otimes \cA$ mediates between
the MC coupling $(V, \mu)$ and the $\cA^!$-module
$(V, \rho)$ via $\mu = (\rho \otimes \mathrm{id})(\tau)$.
The quasi-equivalence follows from the bar-cobar Quillen
equivalence applied to modules.
\end{proof}

\begin{remark}[Off the Koszul locus]
\label{rem:off-koszul-locus}
For non-Koszul algebras \textup{(}e.g.\ simple quotients
$L_k(\fg)$ at admissible levels\textup{)}, the MC coupling category
$\operatorname{Line}(\cA)$ is still well-defined and is still the
correct primitive. What fails is the algebraic presentation as
$\Perf(\cA^!)$: the Koszul dual may not exist, or it may exist
but the comparison functor may not be an equivalence.
The Koszulness characterization programme
\textup{(}$10{+}1{+}1$ conditions\textup{)} is precisely the
classification of when the algebraic presentation works.
\end{remark}

% -------------------------------------------------------------------
\subsection{From local collars to global factorization}
\label{subsec:local-to-global-factorization}
% -------------------------------------------------------------------

\begin{construction}[Local collar categories]
\label{constr:local-collar-categories}
\index{collar category|textbf}
Let $(X, D, \tau)$ be a tangential log curve and
$I_p \cong \R$ the boundary interval at a puncture~$p$.
For a small subinterval $J \subset I_p$, the \emph{collar category}
$\cC_J := \operatorname{Line}(\cA_J)$
is the MC coupling category of the restriction of the bulk algebra
to the formal neighborhood of~$J$.
\end{construction}

\begin{theorem}[$E_1$-factorization on disjoint intervals;
 \ClaimStatusProvedHere]
\label{thm:e1-factorization-disjoint}
For disjoint subintervals $J_1, J_2 \subset I_p$ with
$\sup J_1 < \inf J_2$, the external tensor product
\[
 \cC_{J_1} \boxtimes \cC_{J_2}
 \;\xrightarrow{\;\sim\;}
 \cC_{J_1 \cup J_2}
\]
is an equivalence. This is the $E_1$-factorization structure of
the open color: the ordered interval composition on~$\R$ gives
the associative monoidal structure on the line-operator category.
\end{theorem}

\begin{proof}
The interval $J_1 \cup J_2$ (with a gap between $J_1$ and $J_2$)
retracts onto $J_1 \sqcup J_2$ (disjoint union). The MC
coupling category is homotopy-invariant
(Proposition~\ref{prop:mc-homotopy-invariance}), so
$\cC_{J_1 \cup J_2} \simeq \cC_{J_1 \sqcup J_2}
= \cC_{J_1} \boxtimes \cC_{J_2}$.
\end{proof}

\begin{theorem}[Extension across collisions;
 \ClaimStatusProvedHere{} on the chirally Koszul locus]
\label{thm:collision-extension}
On the chirally Koszul locus, the meromorphic tensor product
\[
 \ell \otimes_z \ell'
 \;:=\;
 (V_\ell \otimes V_{\ell'},\;
 \mu_\ell(z) + \mu_{\ell'}(0) + r_{\ell,\ell'}(z))
\]
extends the $E_1$-factorization structure across the collision
strata of the ordered boundary configuration space, where
$r_{\ell,\ell'}(z)$ is the universal quantum correction computed
by tree-level ladder diagrams. The extension is associative:
\[
 (\ell_1 \otimes_{z_1} \ell_2) \otimes_{z_2} \ell_3
 \;\simeq\;
 \ell_1 \otimes_{z_1 + z_2} (\ell_2 \otimes_{z_2} \ell_3).
\]
\end{theorem}

\begin{proof}
The correction $r_{\ell,\ell'}(z)$ is meromorphic in~$z$ by the
one-loop exactness property of the perturbative 3d~HT theory
(Theorem~\ref{thm:one-loop-koszul}). Associativity follows from
the Yang--Baxter equation for the spectral $R$-matrix
$R(z) = 1 + r(z)/z + \cdots$, which is proved by Stokes' theorem
on $\FM_3(\C)$
(Theorem~\ref{thm:spectral-ybe}).
\end{proof}

\begin{conjecture}[Global Weiss descent]
\label{conj:weiss-descent}
\index{Weiss descent!boundary factorization|textbf}
The local collar categories $\{\cC_J\}_{J \subset I_p}$,
equipped with the $E_1$-factorization on disjoint intervals and the
meromorphic tensor extension across collisions, satisfy Weiss
descent on the boundary configuration space: the natural map
\[
 \cC_{\mathrm{global}}
 \;\longrightarrow\;
 \holim_{J \in \mathrm{Weiss}(I_p)} \cC_J
\]
is an equivalence. This is the categorical boundary analogue of
the Beilinson--Drinfeld factorization axiom for algebras.
\end{conjecture}

\begin{remark}[What is proved vs.\ what is conjectured]
\label{rem:local-vs-global-honest}
The local theory is a theorem:
\begin{itemize}[nosep]
\item Collar categories are well-defined
 (Definition~\ref{def:mc-coupling-category}).
\item $E_1$-factorization on disjoint intervals is proved
 (Theorem~\ref{thm:e1-factorization-disjoint}).
\item Extension across collisions is proved on the Koszul locus
 (Theorem~\ref{thm:collision-extension}).
\item The Koszul comparison $\operatorname{Line}(\cA)
 \simeq \Perf(\cA^!)$ is proved on the Koszul locus
 (Theorem~\ref{thm:koszul-comparison-mc}).
\end{itemize}
The global theory is a conjecture:
\begin{itemize}[nosep]
\item Weiss descent on the boundary Ran space
 (Conjecture~\ref{conj:weiss-descent}).
\item The categorical Ran-space formalism in the boundary setting
 (requires the Raskin--Beraldo framework adapted to
 real-oriented blowups).
\end{itemize}
The honest description: the primitive object is defined locally
as a theorem and is expected to globalize by general categorical
arguments that have not yet been carried out in the boundary
setting.
\end{remark}

% -------------------------------------------------------------------
\subsection{The Kaluza--Klein transfer and non-planar corrections}
\label{subsec:kk-transfer-nonplanar}
% -------------------------------------------------------------------

\begin{theorem}[Tree-level KK transfer;
 \ClaimStatusProvedElsewhere]
\label{thm:tree-level-kk}
\index{Kaluza--Klein transfer!tree-level|textbf}
The Kaluza--Klein reduction of $5$d holomorphic Chern--Simons on
$\C \times S^3$ to a $3$d HT theory on $\C \times \R$ is
controlled at tree level by the $A_\infty$ structure on
$H^{b_0, \bullet}(S^3)$ obtained by homotopy transfer from the dg
algebra $\Omega^{0,\bullet}(S^3)$. The transferred products
$m_2^{\mathrm{tree}}, m_3^{\mathrm{tree}}, \ldots$ are computed by
tree-level Feynman diagrams on~$S^3$.
\end{theorem}

\begin{conjecture}[Non-planar KK transfer]
\label{conj:nonplanar-kk}
\index{Kaluza--Klein transfer!non-planar|textbf}
The full non-planar structure \textup{(}at all loop orders in
$1/N$\textup{)} is controlled by a \emph{deformed} $A_\infty$
structure on $H^{b_0, \bullet}(S^3)$, where the deformation
parameter is $\varepsilon_1 \sim 1/N$
\textup{(}the $\Omega$-background parameter\textup{)}. The
deformed operations $m_k^{(g)}$ at genus~$g$ are the $g$-loop
corrections to the KK transfer. The identification of these
deformed operations with the known non-planar corrections to the
twisted holographic OPE is a checkable prediction of the
programme.
\end{conjecture}

\begin{remark}[Scope]
\label{rem:kk-transfer-scope}
Theorem~\ref{thm:tree-level-kk} is standard
\textup{(}Costello--Li, Costello--Gaiotto\textup{)}.
Conjecture~\ref{conj:nonplanar-kk} is the programme's own
construction: it uses the general theory of $A_\infty$ algebras
in families to produce the deformed structure, but the
identification of the deformation with the physical non-planar
corrections has not been verified by explicit computation at any
loop order beyond tree level.
\end{remark}
